\section{Limitations}
\label{sec:limitations}

This is a small pilot: \NumPrimaryFaces{} faces support the main comparisons,
each model--control cell has fewer than ten faces, and the chained probe uses one
additional face. Multiple outputs from one face are correlated; the clustered
interval addresses that dependence descriptively but cannot provide broad
population inference. Each cell contains one generated output, and two endpoints
ignore the supplied seed, so stochastic variation is unknown.

The localization ratio measures CIELAB pixel change, not anatomical direction,
clinical plausibility, realism, or expected postoperative appearance. Its keep
zone covers only part of the face and it abstains on profile views, precisely where
rhinoplasty visualization is important. Compositing guarantees off-mask
preservation, making that result a verification of system behavior rather than a
test of model quality. ArcFace's 0.6 line is heuristic, and the postoperative
cosines lack an input baseline; neither is a validated surgical-outcome metric.

No surgeon ratings were collected. Clinical plausibility therefore rests neither
on expert judgment nor on a clinically validated automated score. The demographic
composition of the face set was not recorded, so the results cannot be generalized
across skin tones, ages, or sexes; commercial face pipelines have documented
demographic disparities~\cite{buolamwini2018gendershades}. Public visibility of
the source photographs does not itself establish republication, hosted-processing,
or research-use permission, which must be resolved before submission.

Finally, hosted endpoints and aliases can change without notice. Run dates and
endpoint identifiers aid auditing but do not pin every underlying model revision.
The public repository contains the canonical score table and paper analysis, not
the complete generation/scoring implementation or source photographs, so the
study is not presently reproducible end to end.
